% Options for packages loaded elsewhere
\PassOptionsToPackage{unicode}{hyperref}
\PassOptionsToPackage{hyphens}{url}
\documentclass[
  11pt,
  a4paper,
]{article}
\usepackage{xcolor}
\usepackage[margin=25mm]{geometry}
\usepackage{amsmath,amssymb}
\setcounter{secnumdepth}{-\maxdimen} % remove section numbering
\usepackage{iftex}
\ifPDFTeX
  \usepackage[T1]{fontenc}
  \usepackage[utf8]{inputenc}
  \usepackage{textcomp} % provide euro and other symbols
\else % if luatex or xetex
  \usepackage{unicode-math} % this also loads fontspec
  \defaultfontfeatures{Scale=MatchLowercase}
  \defaultfontfeatures[\rmfamily]{Ligatures=TeX,Scale=1}
\fi
\usepackage{lmodern}
\ifPDFTeX\else
  % xetex/luatex font selection
\fi
% Use upquote if available, for straight quotes in verbatim environments
\IfFileExists{upquote.sty}{\usepackage{upquote}}{}
\IfFileExists{microtype.sty}{% use microtype if available
  \usepackage[]{microtype}
  \UseMicrotypeSet[protrusion]{basicmath} % disable protrusion for tt fonts
}{}
\makeatletter
\@ifundefined{KOMAClassName}{% if non-KOMA class
  \IfFileExists{parskip.sty}{%
    \usepackage{parskip}
  }{% else
    \setlength{\parindent}{0pt}
    \setlength{\parskip}{6pt plus 2pt minus 1pt}}
}{% if KOMA class
  \KOMAoptions{parskip=half}}
\makeatother
\usepackage{longtable,booktabs,array}
\usepackage{caption}
\captionsetup[table]{skip=6pt}
\newcounter{none} % for unnumbered tables
\usepackage{calc} % for calculating minipage widths
% Correct order of tables after \paragraph or \subparagraph
\usepackage{etoolbox}
\makeatletter
\patchcmd\longtable{\par}{\if@noskipsec\mbox{}\fi\par}{}{}
\makeatother
% Allow footnotes in longtable head/foot
\IfFileExists{footnotehyper.sty}{\usepackage{footnotehyper}}{\usepackage{footnote}}
\makesavenoteenv{longtable}
\usepackage{graphicx}
\makeatletter
\newsavebox\pandoc@box
\newcommand*\pandocbounded[1]{% scales image to fit in text height/width
  \sbox\pandoc@box{#1}%
  \Gscale@div\@tempa{\textheight}{\dimexpr\ht\pandoc@box+\dp\pandoc@box\relax}%
  \Gscale@div\@tempb{\linewidth}{\wd\pandoc@box}%
  \ifdim\@tempb\p@<\@tempa\p@\let\@tempa\@tempb\fi% select the smaller of both
  \ifdim\@tempa\p@<\p@\scalebox{\@tempa}{\usebox\pandoc@box}%
  \else\usebox{\pandoc@box}%
  \fi%
}
% Set default figure placement to htbp
\def\fps@figure{htbp}
\makeatother
\setlength{\emergencystretch}{3em} % prevent overfull lines
\providecommand{\tightlist}{%
  \setlength{\itemsep}{0pt}\setlength{\parskip}{0pt}}
\usepackage{amsmath,amssymb}
\usepackage{xurl}
\usepackage{fvextra}
\DefineVerbatimEnvironment{verbatim}{Verbatim}{breaklines=true,fontsize=\small}
\urlstyle{same}
\setlength{\emergencystretch}{2em}
\AtBeginDocument{\hypersetup{colorlinks=true,linkcolor=blue,urlcolor=blue}}

% Compatibility: Pandoc image alt metadata with the retained TeX Live 2020.
\makeatletter
\define@key{Gin}{alt}{}
\makeatother
\usepackage{bookmark}
\IfFileExists{xurl.sty}{\usepackage{xurl}}{} % add URL line breaks if available
\urlstyle{same}
% fallback for those not using the hyperref driver hyperxmp:
\makeatletter
\@ifundefined{xmpquote}{\newcommand{\xmpquote}[1]{#1}}{}
\makeatother
\hypersetup{
  pdftitle={Minimum Widths for Unit Squares in Thin Rectangles},
  pdfauthor={Hiroaki Hosono (working draft)},
  hidelinks,
  pdfcreator={LaTeX via pandoc}}

\title{Minimum Widths for Unit Squares in Thin Rectangles}
\author{Hiroaki Hosono (working draft)}
\date{25 October 2026}

\begin{document}
\maketitle

\section{Abstract}\label{abstract}

We present a computer-assisted determination of the minimum width needed
to pack \(N\) unit squares, with arbitrary rotations and pairwise
disjoint interiors, in a rectangle of height \(1851/1000\le a<2\). For
odd \(N\ge11\), the width is \(4+2(N-2)\sqrt2-(N-1)a\). For even
\(N\ge12\), it is the minimum of an explicit rationally parametrised
function, attained by a staggered configuration with two upright end
caps. The representative minimising angle is unique. The proof separates
a finite certificate restricting every candidate packing to a specified
configuration class from analytical minimisation and attainment. We
describe the exact acceptance obligations and the retained two-system
evidence.

\section{1. Introduction and notation}\label{introduction-and-notation}

Let \(\nu(x,a)\) be the maximum number of closed unit squares contained
in \([0,x]\times[0,a]\) with pairwise disjoint interiors, allowing
arbitrary orientations. Contacts are permitted. Write

\[
B_N(a)=\inf\{x>0:\nu(x,a)\ge N\}.
\]

Friedman's rectangle problem {[}1{]} records constructions and boundary
pieces for small numbers of squares. Related counting bounds for unit
squares in rectangles are studied in {[}2,3{]}. The lower bound in the
present paper uses the finite verification and analytical argument
described below.

Our question fixes a height strictly below two and determines the width
uniformly in the number of squares.

No priority claim is made; a broader literature search remains future
work.

\section{2. Main theorem}\label{main-theorem}

Measure \(\theta\) counterclockwise from the positive horizontal axis to
a side direction of a unit square. Orientations differing by \(\pi/2\)
describe the same square. Represent an orientation by
\(t=\tan(\theta/2)\) and let

\[
\mathcal O=\left[2\arctan\frac{4339}{12800},\;2\arctan\frac{4553}{8800}\right].
\]

For \(\theta\in\mathcal O\), put \(c=\cos\theta\), \(s=\sin\theta\),
\(h=c+s\), \(b=3a/2-1\), and

\[
P=\frac{2h-a}{h^2-1},\qquad
Q=1+\frac{2h-b}{h^2-1},\qquad
A=\frac{2-a}{4}(\tan\theta-\cot\theta).
\]

Define

\[
G_N(\theta;a)=(N-4)P+2Q-2|A|.
\]

All denominators are positive on this angle interval. The symbol \(h\)
denotes a square's bounding-box side length, whereas \(a\) is the
container height.

\textbf{Theorem 2.1 (exact thin-strip width; computer-assisted).} For
every \(1851/1000\le a<2\):

\begin{enumerate}
\def\labelenumi{\arabic{enumi}.}
\tightlist
\item
  For odd \(N\ge11\), \[
  B_N(a)=4+2(N-2)\sqrt2-(N-1)a.
  \]
\item
  For even \(N\ge12\), \[
  B_N(a)=\min_{\theta\in\mathcal O}G_N(\theta;a).
  \] The representative minimiser \(\theta_N(a)>\pi/4\) is unique and
  lies below \(2\arctan(43/100)\). Its reflection \(\pi/2-\theta_N(a)\)
  gives the other minimising orientation. The value is attained by a
  staggered packing.
\end{enumerate}

Both formulas are attained by two upright end caps and \(N-2\) interior
squares at one common angle, placed on alternating walls. For odd \(N\),
the two caps touch the same wall and the optimal common angle is
\(\pi/4\). For even \(N\), the caps touch opposite walls and the optimal
angle is selected by the one-variable minimum in the theorem. Figure 1
shows examples of this general construction; the proof of attainment for
every size is given in Section 5.

An interactive illustration in the online version allows the number of
squares and the height to vary. Decimal widths and angles in that
illustration are floating-point approximations; the exact claims are the
formulas above. The finite range of the controls does not limit the
theorem. The illustration displays the attaining family and is not a
proof computation.

The theorem combines the retained two-system results on
\([1851/1000,19/10]\) and \([19/10,2)\), including the separate coverage
for \(N=11\). The mathematical argument is given below, with its finite
proof obligations exposed in Section 6.

\section{3. Certified restriction to a configuration
class}\label{certified-restriction-to-a-configuration-class}

First order the squares by nondecreasing horizontal centre coordinate;
ties may be ordered arbitrarily. Use square symmetry to represent each
angle by \(0\le\theta\le\pi/2\), with \(t=\tan(\theta/2)\). In this
representation, \textbf{near-upright} means

\[
t\in\widehat{\mathcal U}=[0,7/256]\cup[125/132,1].
\]

For a square with centre height \(y\) and bounding-box side
\(h_\theta=\cos\theta+\sin\theta\), define its normalised height

\[
\rho=\frac{y-h_\theta/2}{a-h_\theta}\in[0,1].
\]

The denominator is positive in the theorem range. The closed lower and
upper height labels are \(\rho\in[0,1/2]\) and \(\rho\in[1/2,1]\). A
midpoint may receive either label in the closed cover. A
\textbf{sequence in the configuration class} (a \emph{type sequence} in
the source records) has the half-angle parameters of its two end squares
in \(\widehat{\mathcal U}\), every interior angle in \(\mathcal O\), and
alternating height labels along all consecutive pairs, including the
pairs at the caps. The class restriction is the assertion that every
packing at or below the threshold stated next has such a type sequence.

Put \(g=2\sqrt2-a\) and \(F_0=4-2\sqrt2\). The retained certificate
establishes the following assertion for every \(N\ge11\): a packing of
width at most

\[
(N-1)g+F_0
\]

has two near-upright end caps, its \(N-2\) interior squares have angles
in \(\mathcal O\), and the centres of the interior squares alternate
between the closed lower and upper halves of the container (the midpoint
belongs to both). Every sequence outside the class has strictly greater
width.

This is an all-sequence assertion. The verification covers angle/height
cells and all transitions, not just configurations found in a search.
The separating-axis theorem (SAT), which tests disjointness by
projections onto edge normals, gives exact bounds for transitions
outside the class. Rational potentials, found by linear programming (LP)
and checked in exact arithmetic, redistribute these bounds along a
sequence. Calibrated transitions within the class use angular variation.
A saturated dynamic-programming (DP) layer and its checked fixed point
cover arbitrarily large \(N\).

The extension from finitely many sizes to every larger size is an
induction, not an inference from sampled values. After the chosen edge
layer, consider the finite augmented state graph, including the flags
and previous-edge information needed by the transition rules. Let
\(b_z\) be the bound at entry to the saturated layer, \(w_{zz'}\) a
verified transition cost, and \(V_z\) the final saturated bound. The
exact fixed-point check establishes

\[
V_z\le b_z,\qquad V_{z'}\le V_z+w_{zz'}
\quad\text{for every legal transition }z\to z'.
\]

At entry the accumulated path cost is at least \(V_z\). If that holds at
one state, adding the next verified transition cost proves it at the
next state. Thus it holds after any number of further transitions.
Adding each legal terminal bound gives the required residual-width bound
for arbitrarily long sequences outside the class. The start and end
conventions convert edge counts to square counts; their two end rules
are why a saturated edge layer must not be labelled as a finite count of
squares. The finite sizes before this tail are checked separately, as
listed in Section 6. Every legal transition, including entry into a run
of compatible interior states, must be present for this induction to
apply. This argument uses the checked finite vector inequalities, and
does not assume that all chord-transition costs are nonnegative.

The threshold is sufficient for the equality theorem as follows. The
\(45^\circ\) construction has width \((N-1)g+F_0\). For even \(N\),
\(G_N(\pi/4;a)\) equals that threshold, so \(\min G_N\) is no larger. An
arbitrary packing wider than the threshold already satisfies the claimed
lower bound. Any packing at or below the threshold is restricted to this
class, where Section 4 applies. Section 5 constructs a packing attaining
the resulting value. Thus the class restriction alone is not being
substituted for minimisation or attainment.

The finite verification includes transitions into runs of compatible
interior states as well as transitions within those runs. Near height
two it uses a corrected absolute-value inequality for nearly upright
squares. Earlier computations with the superseded signed inequality are
excluded from the evidence.

\section{4. The lower bound within the configuration
class}\label{the-lower-bound-within-the-configuration-class}

\subsection{4.1. Pair and cap
inequalities}\label{pair-and-cap-inequalities}

For interior angles \(\theta_1,\ldots,\theta_{N-2}\) define

\[
\mathrm{TV}=\sum_{i=1}^{N-3}|\theta_{i+1}-\theta_i|.
\]

The retained pair calibration gives a lower bound per radian of
\(\theta\). Its constant can be chosen as

\[
k=\frac15\quad\left(\frac{1851}{1000}\le a\le\frac{19}{10}\right),
\qquad
k=\frac{227}{2000}\quad\left(\frac{19}{10}\le a<2\right).
\]

Here is the summed inequality explicitly. Put \(r=N-2\) and label the
interiors from left to right. The cap centres have horizontal
coordinates \(x_0\) and \(x_{r+1}\). Let \(\sigma_0=+1\) when the first
cap-to-interior step goes from a lower-labelled square to an
upper-labelled one, and \(\sigma_0=-1\) for the reverse step. Set

\[
\sigma_i=(-1)^i\sigma_0\quad(0\le i\le r),\qquad
d_i=x_{i+1}-x_i.
\]

The gap functions \(f_\pm\) and \(c_\pm\) are defined explicitly in
Section 5. The calibrated interior-pair inequality is

\[
d_i\ge\frac{f_{\sigma_i}(\theta_i)+f_{\sigma_i}(\theta_{i+1})}{2}
+k|\theta_{i+1}-\theta_i|\quad(1\le i<r).
\]

Thus calibration means a uniform strengthening of the common-angle gap
bound by a term proportional to angle difference. If \(h_L,h_R\) are the
bounding-box side lengths of the original two caps, the cap comparison
gives

\[
\frac{h_L}{2}+d_0\ge\frac12+c_{\sigma_0}(\theta_1),\qquad
d_r+\frac{h_R}{2}\ge c_{\sigma_r}(\theta_r)+\frac12.
\]

The identities

\[
P=\frac{f_++f_-}{2},\qquad
T_\sigma=\frac12+c_\sigma+\frac12f_{-\sigma}=Q+\sigma A
\]

then give the occupied-width bound

\[
W\ge\sum_{j=2}^{r-1}P(\theta_j)
+Q(\theta_1)+Q(\theta_r)
+\sigma_0 A(\theta_1)+\sigma_r A(\theta_r)+k\,\mathrm{TV}.
\]

Because \(\sigma_r=(-1)^r\sigma_0\), the endpoint signs agree for even
\(N\) and are opposite for odd \(N\). This is the sign rule used in the
next subsection. The same certified bounds give
\(|\partial_\theta A|\le k\).

The cap comparison is applied to pairs at their original positions. It
does not assume that simultaneously moving arbitrary caps preserves a
packing. The lower-height cap comparison is a separate certified
dependency; the upper-height argument must not be extrapolated below 1.9
without it.

\subsection{4.2. Reduction to a common
angle}\label{reduction-to-a-common-angle}

Choose an interior angle \(\bar\theta\) closest to \(\pi/4\). Both \(P\)
and \(Q\) decrease as \(h\) increases. Indeed,

\[
\frac{d}{dh}\frac{2h-z}{h^2-1}
=-\frac{2(h^2-zh+1)}{(h^2-1)^2}<0
\]

for \(z=a\) or \(z=b\), since \(z<2\) in the claimed range. Consequently
each \(P\)- and \(Q\)-term is at least its value at \(\bar\theta\).

For even \(N\), the two endpoint \(A\)-terms have the same sign pattern.
Replacing them by \(-2|A(\bar\theta)|\) loses at most

\[
k\bigl(|\theta_1-\bar\theta|+|\theta_{N-2}-\bar\theta|\bigr)
\le k\,\mathrm{TV}.
\]

The variation contribution absorbs this loss. The occupied width is
therefore at least \(G_N(\bar\theta;a)\), and hence at least its
minimum.

For odd \(N\), the endpoint \(A\)-terms have opposite signs. Their
difference is absorbed by \(k\,\mathrm{TV}\). The remainder is minimised
when \(h=\sqrt2\), giving

\[
(N-4)P(\pi/4)+2Q(\pi/4)
=4+2(N-2)\sqrt2-(N-1)a.
\]

\subsection{4.3. Uniqueness and the critical
equation}\label{uniqueness-and-the-critical-equation}

In this paragraph differentiation is with respect to \textbf{\(t\)}, not
per radian of \(\theta\). On

\[
I=[\sqrt2-1,4553/8800],
\]

the retained exact polynomial checks establish

\[
P_{tt}>0,\qquad P_{tt}+2Q_{tt}-2A_{tt}>0.
\]

On the open upper-angle side \(t>\sqrt2-1\), \(A\ge0\), and

\[
G_{N,tt}=(N-5)P_{tt}+(P_{tt}+2Q_{tt}-2A_{tt})>0.
\]

The full function has an absolute-value cusp at \(t_0=\sqrt2-1\). The
derivative here is the right derivative of the upper-angle branch:

\[
G_{N,t}(t_0+;a)=-(2-a)(2+\sqrt2)<0,
\]

while the checked derivative at \(t=43/100\) is positive for \(N\ge12\).
This yields a unique representative critical point below that endpoint.
The reflection of the lower-angle part of \(\mathcal O\) lies in the
checked upper interval, so it adds no smaller value. Reflection gives
the other minimising orientation.

The rational parametrisation of \(c,s\) turns the critical-point
equation into a polynomial equation in \(t\), with rational polynomial
coefficients in \(a\), after multiplication by positive denominators.
The minimising parameter is therefore an algebraic function of \(a\).
When \(a\) is algebraic, its parameter and minimum width are algebraic
numbers. For an arbitrary real \(a\), no algebraicity over the rationals
is asserted, and the angle \(2\arctan t\) is not being called an
algebraic number. Convexity and endpoint signs here are certificate
dependencies, not deductions from a plot.

\section{5. Attaining configurations}\label{attaining-configurations}

Use two upright caps and \(N-2\) interiors at one angle, on alternating
walls. Let

\[
f_+=\frac{1-s(a-h)}c,\qquad
f_-=\frac{1-c(a-h)}s,
\]

\[
c_+=\frac{(1+h)/2-s(a-(1+h)/2)}c,\qquad
c_-=\frac{(1+h)/2-c(a-(1+h)/2)}s.
\]

Interior centres have height \(h/2\) or \(a-h/2\). Each cap has centre
height \(1/2\) or \(a-1/2\), on the wall opposite its neighbour.
Starting with cap centre \(x=1/2\), take consecutive horizontal
increments \(c_\pm\) at the caps and \(f_\pm\) between interiors,
according to the alternating wall pattern.

Adjacent pairs touch on an appropriate face normal. On the angle
interval containing every minimiser, the retained attainment checks
establish

\[
f_\pm>h/2,\qquad
c_\sigma+f_{\sigma'}>(1+h)/2,\qquad
c_\sigma>(h-1)/2
\]

for the required sign choices. The first two inequalities separate
nonadjacent pairs horizontally; longer chains add positive increments.
The last makes the caps outermost. The cap--cap pair is also separated,
since a cap increment and two interior increments already exceed one.

Containment follows from \(h\le\sqrt2<a\). Thus the width is exactly one
plus the sum of the horizontal increments. For even \(N\), choose the
starting wall to obtain

\[
1+2c_{\sigma_0}+f_{-\sigma_0}
 +(N-4)(f_++f_-)/2
=(N-4)P+2Q-2|A|.
\]

For odd \(N\), use \(\theta=\pi/4\); the same sum is the formula in
Theorem 2.1. Choosing the critical orientation in the even case gives
exact attainment. The alternating sequence puts the two upright caps on
the same wall for odd numbers of squares and on opposite walls for even
numbers (Figure 1).

For fixed \(t\), the individual expressions in the analysis and
attainment checks are affine in \(a\). The accepted endpoint
certificates therefore extend pointwise through the two height
intervals. This is applied to individual inequalities, not to a presumed
endpoint extremum of an arbitrary minimum of affine branches.

A nearby rational-half-angle placement certificate may have a small
positive width slack. Such a certificate is evidence of a packing near
the optimum, not equality of its width with the critical value.

\par\medskip
\noindent\begin{minipage}{\linewidth}
\centering
\pandocbounded{\includegraphics[keepaspectratio,alt={Alternating-wall constructions for odd and even numbers of squares.}]{figures/parity-packing.pdf}}

\par\smallskip\raggedright
\emph{Figure 1. The construction of Section 5 at height 1.9: eleven
squares at 45 degrees (the odd optimum), and twelve squares at the
illustrative rational half-angle parameter 427/1000. Orange squares are
the two upright end caps; blue squares alternate between the walls. The
even panel illustrates a feasible construction at the displayed angle,
not the exact critical-angle optimum. Coordinates are rounded only for
drawing; the figure is not a packing certificate.}
\end{minipage}\par\medskip

\par\medskip
\noindent\begin{minipage}{\linewidth}
\centering
\pandocbounded{\includegraphics[keepaspectratio,alt={Height-width formulas for eleven squares and illustrative twelve-square constructions.}]{figures/height-width.pdf}}

\par\smallskip\raggedright
\emph{Figure 2. Height and width on the theorem interval. The odd curve
is the exact formula for eleven squares. The even curve is the width of
the twelve-square construction at the fixed half-angle parameter
427/1000; it is an upper bound, while the theorem takes the minimum over
all admissible angles. The right endpoint at height two is excluded.
Lines are drawn from the stated formulas and do not establish the finite
lower-bound computation.}
\end{minipage}\par\medskip

\section{6. Exact computations and
reproducibility}\label{exact-computations-and-reproducibility}

Appendix A lists the requirements for independent reproduction. The
mathematical coverage and the obligations for a complete verification
are set out here.

\subsection{6.1. Verification coverage}\label{verification-coverage}

The lower-height verification records document checks by a first
implementation and a specification-derived second implementation,
followed by mathematical review. It reports complete interval and Sturm
checks separately from the 36 finite placement tests. Those placements
are not a proof of all heights or all \(N\).

For \(a\in[1851/1000,1852/1000]\), the combined finite-layer and tail
verification establishes \(N\ge24\) using its individual finite layers
together with its saturated tail; the first-system output labels that
tail \(N\ge26\), rather than \(N\ge24\). Separate fixed-\(N\) checks
cover \(11\le N\le26\). Their union covers every \(N\ge11\). On
\([1852/1000,1854/1000]\), the chord runs cover \(N\ge12\) and separate
runs cover \(N=11\). Later-height runs cover the remainder. The
corrected absolute-form final interval and the separate \(N=11\)
subdivision at height 1.9 must be included.

The verification records document that the second implementation was
written from the specifications without using the first implementation's
code.

\subsection{6.2. Acceptance obligations}\label{acceptance-obligations}

A complete replay must check the closed state cover, all SAT branches,
potential coverage, every transition (including entry to and exit from
runs of compatible interior states), the integer rounding model, and the
saturated-layer fixed point. It must also check the full-domain
calibration and cap inequalities, the polynomial signs and root counts,
the affinity argument and the attainment inequalities. Positive output
on selected intervals, floating LP feasibility, placement samples or a
root count alone are insufficient.

\subsection{6.3. Availability of the computational
supplement}\label{availability-of-the-computational-supplement}

The implementations use Python 3, SymPy and NumPy. A self-contained
computational supplement must specify their versions, the complete
source code, and the identities of every accepted input and output. That
version-pinned supplement is not yet available with this paper. The
public paper and its TeX bundle provide the exposition and figures. The
complete verification-code and input-data archive remains to be
released. Appendix A states what a complete archive must contain and
what each verification must establish.

\section{7. Scope and computational
supplement}\label{scope-and-computational-supplement}

The theorem covers every fixed height \(1851/1000\le a<2\) and every
integer \(N\ge11\), with the odd and even formulas stated in Theorem
2.1. The left height endpoint is included and height two is excluded.
Other heights and smaller values of \(N\) are outside the statement.

This is a computer-assisted working paper and has not been externally
peer reviewed. The assembled text has received a limited independent
consistency review against the verification records. A self-contained,
version-pinned computational supplement and a complete replay inventory
remain to be released. Preparing the text and illustrations did not
include a new execution of the proof computations.

\section{8. References}\label{references}

{[}1{]} Erich Friedman, \emph{Problem of the Month (June 2010)}, Math
Magic. \url{https://erich-friedman.github.io/mathmagic/0610.html}.

{[}2{]} Hiroshi Nagamochi, \emph{Packing Unit Squares in a Rectangle},
The Electronic Journal of Combinatorics \textbf{12} (2005), Research
Paper R37. \url{https://doi.org/10.37236/1934}.

{[}3{]} Hakan Karakuş, \emph{A counterexample to Nagamochi's scoring
lemma and a new rectangle packing bound}, arXiv:2609.37410v1 (2026).
\url{https://arxiv.org/abs/2609.37410v1}.

\fontsize{9.5}{10.8}\selectfont
\section{Appendix A. Requirements for independent
reproduction}\label{appendix-a.-requirements-for-independent-reproduction}

\subsection{A.1. Verification
obligations}\label{a.1.-verification-obligations}

A complete computational supplement must connect each proof ingredient
to a precisely specified verification and an unambiguous set of accepted
input and output data. The following obligations describe the
mathematical content; the complete archive has not yet been released.

{\def\LTcaptype{none} % do not increment counter
\begin{longtable}[]{@{}
  >{\raggedright\arraybackslash}p{(\linewidth - 2\tabcolsep) * \real{0.5000}}
  >{\raggedright\arraybackslash}p{(\linewidth - 2\tabcolsep) * \real{0.5000}}@{}}
\toprule\noalign{}
\begin{minipage}[b]{\linewidth}\raggedright
Ingredient
\end{minipage} & \begin{minipage}[b]{\linewidth}\raggedright
Required verification
\end{minipage} \\
\midrule\noalign{}
\endhead
\bottomrule\noalign{}
\endlastfoot
state cover, exact bounds, DP and saturation rules & prove closed
coverage, sound rounding and the fixed-vector inequalities \\
upper-height analysis and attainment & verify polynomial signs and
nonintersection throughout the stated intervals \\
upper-height coverage, including the final interval and \(N=11\) &
record both the exceptional finite cases and the arbitrarily large
tail \\
lower-height configuration-class restriction & check every legal
transition, including entries into interior runs \\
lower-height calibration and cap comparison & verify the quantitative
pair inequalities and cap bounds at original positions \\
lower-height analysis and attainment & verify the angle interval,
affine-height extension and all required separation bounds \\
upper-height calibration & verify the common-angle and angle-variation
bounds uniformly \\
\end{longtable}
}

\subsection{A.2. Archive and replay
requirements}\label{a.2.-archive-and-replay-requirements}

The released archive must specify the full source code and dependency
versions, the exact state covers and potential data, all accepted inputs
and outputs, the height intervals and finite-count ranges, and the
checked saturation conditions. It must document how the two
implementations were developed independently and identify which output
establishes each proof obligation.

Reproduction must cover all separating-axis branches, all legal
transitions and terminal rules, exact integer rounding, every
fixed-point inequality, the analytical polynomial checks, and the
attaining configurations. A selection of examples or floating-point
optimization results does not replace these checks.

In particular, fixed-\(N\) checks do not replace the separately required
saturation argument. The final height interval \([1.999,2)\) requires
the absolute-value form of the near-upright inequality. These
requirements must be explicit in the released verification data.

\end{document}
